\subsection{The Breath of the First Forge, the Spark in the Maker's Hand}
\label{patron:first-forge}

\subsubsection*{Lore}
Before the first hammer struck, before the first wheel turned, before the first seed was planted in the furrow, there was the Breath. It was not flame---flame came later, when the Breath met kindling. It was not knowledge---knowledge came later, when the Breath was shaped into questions. It was the raw, primal urge to \emph{make}: to take what is and turn it into what could be. The caveman breathing on coals to coax fire from ash. The Aeler smith regulating their breath as they hammer, each exhale a rhythm, each inhale a new possibility. The Lethai artisan breathing \emph{anima}} into clay until the vessel remembers it was once alive.

The Breath of the First Forge is the patron of that moment: the spark, the inspiration, the first step that transforms a lump of ore into a blade, a thought into a song, a hope into a plan. It does not care about perfection---perfection is the Clockwork Monad's domain, a slow and, a slow and grinding hunger. The Breath cares about the \emph{act of creation itself}: the joy of building, the surprise of discovery, the stubborn refusal to accept that a thing cannot be done.

It is the patron of:
\begin{itemize}
  \item The smith who tries a new alloy, just to see what happens.
  \item The child who builds a tower of sticks and watches it fall, then builds it again.
  \item The artist who paints the same scene a hundred times, each iteration a different truth.
  \item The parent who whittles a toy for their child, not because it is needed, but because the whittling is a prayer.
\end{itemize}

The Breath is not a god of crafts---that is too narrow. It is the god of \emph{the moment before the craft begins}, the inspiration that drives the hand. It is the breath that warms the coals, the breath that steadies the hammer, the breath that whispers: ``Try.''

\begin{quote}
``Before the first hammer struck, before the first wheel turned, before the first seed was planted in the furrow, there was the Breath. It is not a flame---flame came later, when the Breath met kindling. It is not a plan---plans came later, when the Breath met thought. It is the moment of ignition: the spark, the urge, the refusal to leave the world as it is. It is the patron of everything that has ever been made, and of everything that has ever been made better. But unlike the Clockwork Monad, it does not care about the final form. It cares about the breathing.''
\end{quote}

The Breath is the counterpoint to the Clockwork Monad. Where the Monad seeks perfection through endless refinement, the Breath celebrates the \emph{first attempt}---the flawed, beautiful, human act of creation. A sword made with the Breath's blessing may be imperfect, but it will be \emph{alive} in a way that a Monad-forged blade never can be. It may break, but it will break in service of a story. It may be crude, but it will be loved.

\begin{tcolorbox}[colback=blue!5!white,colframe=blue!70!black,title=The Gray Wanderer's Note]
``The Clockwork Monad asks: 'Can you make it perfect?' The Breath of the First Forge asks: 'Can you make it now?' One is a patient creditor; the other is a hungry friend. Both will cost you---but only one will leave you smiling.''
\end{tcolorbox}

\subsubsection*{Domain Focus}
\begin{itemize}
  \item \textbf{Creative Inspiration:} the spark of new ideas and novel approaches
  \item \textbf{The Joy of Making:} the pleasure and satisfaction of creating
  \item \textbf{Ingenuity:} finding unconventional solutions to problems
  \item \textbf{Primal Craft:} making as a fundamental human (and inhuman) act
  \item \textbf{Anima:} breathing life and spirit into objects
  \item \textbf{Counterpoint to Perfection:} embracing imperfection, iteration, and discovery
\end{itemize}

\subsubsection*{Thiasos and Codex (Runekeeper Options)}
A Runekeeper sworn to the Breath of the First Forge does not seek a perfect familiar. Their \textbf{Thiasos} is a creature that embodies the \emph{spirit of making}---a small, curious, often mischievous presence that delights in creation: a clay fox that reshapes itself each time it is touched, a fire-sprite that dances on the anvil and leaves faint scorch-marks shaped like runes, a living knot of copper wire that unties and reties itself into new patterns, or a small, self-carving wooden bird that sings a different song each dawn. This creature is not bound by precision or efficiency; it is bound by \emph{curiosity}. It requires no fuel, but it needs to be shown new things---new materials, new techniques, new failures. If neglected, it becomes sullen and still (\emph{Compromised}: it refuses to move or sing until you show it something it has never seen before).

The \textbf{Codex} of the Breath is never a finished book. It may be a growing collection of sketches, a scrapbook of failed attempts and happy accidents, a set of interlocking wooden blocks that can be rearranged into infinite configurations, or a roll of birchbark that records only the \emph{first} draft of each Rite (the later refinements are left unwritten). Because the Breath values the moment of ignition, the Codex requires upkeep: the Runekeeper must spend an hour each week attempting to create \emph{something new}---a sketch, a carving, a new combination of materials---or the Rites stored there become static (treat as Compromised until refreshed).

\subsubsection*{Patron's Gift (Imbuement)}
Once per scene as an action, a Runekeeper of the Breath of the First Forge may touch a held item (weapon, tool, or piece of armor) and whisper a word of creation. For the rest of the scene, the item glows faintly with the warmth of a just-lit forge. It grants:
\begin{itemize}
  \item \textbf{+1 die} to any roll involving creative use of the item (improvisation, unconventional applications, or first-time attempts).
  \item The first time each scene the bearer would fail a Craft or Tinker roll, they may reroll it (the item briefly ``suggests'' a new approach).
\end{itemize}
After the scene, the user marks \textbf{+1 Obligation} to the Breath of the First Forge. If they spend the downtime following the scene working on a new creation, they may clear this Obligation for free (the Breath is satisfied by the act of making).

\subsubsection*{Rites of the Breath of the First Forge}
\paragraph{Rite of the Igniting Breath} (Basic, 5 XP) \hfill \\
\emph{Scene; Self/Touch; Yes.} \\
\textbf{Tags:} \texttt{[SPARK]}, \texttt{[INSPIRE]}, \texttt{[HEAT]} \\
\textbf{Materials:} A cold hearth, a dead forge, or a project that has stalled. \\
\textbf{Effect:} Choose one:
\begin{itemize}
  \item \textbf{Spark the Forge:} Relight any non-magical fire or forge as if it had been tended for hours. Gain \textbf{+2 dice} to the next Craft or Tinker roll made at that forge.
  \item \textbf{Inspire the Maker:} A touched ally gains \textbf{+1 die} to their next creative or skill roll and may reroll one failure on that roll. The Breath delights in seeing what they will make.
\end{itemize}
\textbf{Push It:} The spark leaps---the fire catches in an unexpected way. Mark \textbf{1 SB (Diamonds)} as the Breath's enthusiasm spreads beyond your control. Cost: Mark +1 Obligation (total +2) or Mark +1 Fatigue. \\
\textbf{Cost:} Mark +1 Obligation. \\
\textit{Requires:} Familiar (\textit{Invoke:} 1 Boon). \\
\textbf{Invoke:} 1 action. \\
\textbf{Duration:} Scene.

\paragraph{Rite of the Happy Accident} (Basic, 6 XP) \hfill \\
\emph{Instant; Self; Yes.} \\
\textbf{Tags:} \texttt{[CHANCE]}, \texttt{[DISCOVERY]}, \texttt{[BOOST]} \\
\textbf{Materials:} A failed attempt---a cracked pot, a misshapen blade, a smeared sketch. \\
\textbf{Effect:} When you fail a Craft, Tinker, or creative roll, you may invoke this Rite to turn the failure into a discovery. The Breath reveals a hidden lesson in the wreckage:
\begin{itemize}
  \item Gain \textbf{+1 die} on your next attempt at the same task.
  \item Or, gain a \textbf{Boon} (the failure was educational).
  \item Or, learn a minor secret about the materials you were working with (GM grants a hint).
\end{itemize}
\textbf{Push It:} The failure is so spectacular it inspires a completely new approach. Gain \textbf{+2 dice} on your next attempt, but the Breath's enthusiasm marks \textbf{1 SB (Clubs)} as your improvisation draws attention. Cost: Mark +1 Obligation (total +2) or Mark +1 Fatigue. \\
\textbf{Cost:} Mark +1 Obligation. \\
\textit{Requires:} Familiar (\textit{Invoke:} 1 Boon). \\
\textbf{Invoke:} 1 action (reaction to failure). \\
\textbf{Duration:} Instant; Boon or bonus applies to the next attempt.

\paragraph{Rite of the Anima Breath} (Standard, 8 XP) \hfill \\
\emph{Extended; Touch; Yes.} \\
\textbf{Tags:} \texttt{[ANIMA]}, \texttt{[LIFE]}, \texttt{[IMPROVE]} \\
\textbf{Materials:} A creation that is nearly finished, and a willing breath from the maker. \\
\textbf{Effect:} You breathe a spark of life into a creation. Choose one:
\begin{itemize}
  \item \textbf{Temporary Anima:} A mundane object gains a spark of sentience for one scene. It can perform simple tasks (hold a door, fetch a tool, warn of danger) but cannot speak or make complex decisions. Mark \textbf{+1 Obligation}.
  \item \textbf{Permanent Anima:} A creation gains a minor magical property (equivalent to a 2 XP enchantment). This requires a \textbf{Downtime Timer} (4 segments) and the Breath's continued attention. Mark \textbf{+2 Obligation}.
\end{itemize}
\textbf{Push It:} The anima is unusually strong---the object gains a personality and may act independently. Mark \textbf{1 SB (Spades)} as the Breath's gift becomes complicated. Cost: Mark +1 Obligation (total +3) or Mark +1 Fatigue. \\
\textbf{Cost:} Mark +1 Obligation (temporary) or +2 Obligation (permanent). \\
\textit{Requires:} Familiar + Codex (\textit{Invoke:} 1 Boon). \\
\textbf{Invoke:} 1 action (temporary) or downtime action (permanent). \\
\textbf{Duration:} Scene (temporary) or permanent (with upkeep).

\paragraph{Rite of the Unfinished Project} (Standard, 7 XP) \hfill \\
\emph{Scene; Zone; Yes.} \\
\textbf{Tags:} \texttt{[CREATION]}, \texttt{[AREA]}, \texttt{[BOOST]} \\
\textbf{Materials:} A workshop or space littered with half-finished projects. \\
\textbf{Effect:} The Breath blesses a space where making is happening. For the scene:
\begin{itemize}
  \item All Craft, Tinker, and creative rolls in the zone gain \textbf{+1 die}.
  \item Any failure in the zone yields \textbf{+1 Boon} (the Breath loves a learning moment).
  \item Start a \textbf{[4] Inspiration Flood} timer. When it fills, the Breath's enthusiasm becomes overwhelming---everyone in the zone must make a creative roll (any skill) or mark Fatigue from the pressure of \emph{too many ideas}.
\end{itemize}
\textbf{Push It:} The inspiration flood is immediate. Everyone in the zone gains \textbf{+2 dice} to their next creative roll, but they also gain a \textbf{Compulsion: Create} (must attempt at least one creative act before the scene ends). Cost: Mark +1 Obligation (total +3) or Mark +1 Fatigue. \\
\textbf{Cost:} Mark +2 Obligation. \\
\textit{Requires:} Familiar + Codex (\textit{Invoke:} 1 Boon). \\
\textbf{Invoke:} 1 action. \\
\textbf{Duration:} Scene; Timer: 4 segments (Inspiration Flood).

\paragraph{Rite of the Maker's Joy} (Master, 14 XP) \hfill \\
\emph{Extended; Self/Touch; Yes.} \\
\textbf{Tags:} \texttt{[CREATION]}, \texttt{[ANIMA]}, \texttt{[FATE]}, \texttt{[BOOST]} \\
\textbf{Materials:} A creation that you love---one that you made with your own hands and poured your heart into. \\
\textbf{Effect:} You invoke the Breath's deepest gift: the joy of making, made manifest. For one scene:
\begin{itemize}
  \item You gain \textbf{+3 dice} to all Craft, Tinker, and creative rolls.
  \item Your creation (the focus of the Rite) becomes a temporary \textbf{Artifact}---it gains one magical property of your choice (GM approval) for the duration of the scene.
  \item You may reroll any single failure during the scene (once).
\end{itemize}
\textbf{Push It:} The joy becomes contagious. Allies within Near gain \textbf{+1 die} to their creative rolls, but the Breath's attention is on you---mark \textbf{2 SB} for the GM to spend on complications (rival makers, jealous patrons, or the Clockwork Monad taking notice). Cost: Mark +1 Obligation (total +3) or Mark +1 Fatigue. \\
\textbf{Cost:} Mark +2 Obligation. \\
\textit{Requires:} Familiar + Codex + Tier III (\textit{Invoke:} 2 Boons). \\
\textbf{Invoke:} Extended action. \\
\textbf{Duration:} Scene.

\subsubsection*{Cantors \& Cults of the First Forge}
\textbf{The Cantor's Song.} A Cantor of the Breath does not sing melodies; they sing the \emph{sound of making}---the rhythm of a hammer on an anvil, the hiss of hot metal in water, the patient scrape of a chisel, the hum of a spinning wheel. Their voice is percussive, physical, and deeply satisfying to those who work with their hands. They are sought by guilds to bless new workshops and by lone artisans who need to find the joy in a project that has become a chore.

\textbf{The Cult of the First Spark.} Cults of the Breath gather in workshops, forges, and studios. Their rituals are not prayers but \emph{demonstrations}: each member must show something they have made, and the group celebrates the effort, not the quality. A lopsided pot is praised as much as a masterwork sword---what matters is that someone \emph{tried}. Their creed is simple and warm: ``Make something. Make it badly. Make it again.'' Outsiders find them naive; insiders find a rare, safe space to fail and grow. The most devoted members keep a small jar of coals from the forge that first inspired them; they say the coals still warm when they are about to have a good idea.

\textbf{The Cult of the Unfinished} (Hidden). A radical offshoot believes that the Breath's greatest gift is the \emph{infinite potential of the unfinished}. They never complete anything---they leave projects half-done, sketches half-drawn, songs half-composed. They argue that completion is death, and that the Breath lives in the moment between start and finish. They are not malicious, but they can be deeply frustrating to anyone who needs a finished product. They are also, secretly, the Breath's most devoted servants---because they never stop creating.

\subsubsection*{Witchcraft of the First Forge (Lay / Threshold)}
A witch who walks with the Breath is a \emph{spark-keeper}, a mender of the small, quiet moments when a thing becomes possible. Their magic is not flashy---it is the warmth of a hand on a cold tool, the unexpected ease of a difficult cut, the sudden certainty that \emph{this idea will work}. They work with coals, breath, and the dust of failed projects, never with precise measurements or careful plans---because the Breath is not about precision; it is about possibility.

\textbf{The Witch's Tool: The Coal of the First Spark.} A single coal, taken from the forge where the witch first felt the Breath's touch. It never burns down; it simply glows, warm and steady, waiting to be breathed upon. Once per session, the witch may breathe on the coal and speak a single word of creation. For the next scene, that word becomes a \emph{tag} that can be applied to any creative roll: the first use of that tag grants \textbf{+2 dice}, and any failure while it is active reveals a new approach (GM grants a hint). The coal can also be used to rekindle a dying fire or to warm a cold forge.

\textbf{A Hedge Gift: The Moment of Ignition (Basic, 2 XP).} Once per scene, when you are about to begin a creative task (crafting, sketching, composing, even planning), you may declare that you are ``striking the spark.'' Roll the check with \textbf{+1 die}. If you succeed, the Breath blesses your effort---you may reroll one failure on that check.

\subsubsection*{Corruption of the First Forge}
\begin{tabular}{|c|p{6cm}|p{6cm}|}
\hline
\rowcolor{gray!30}\textbf{Tier} & \textbf{Benefit} & \textbf{Cost / Quirk} \\
\hline
1 & \textbf{The Spark of Inspiration:} +1 die to \textit{Tinker/Craft} when starting a new project. & \textbf{Restless Hands:} You must be making something. Sitting idle costs \textbf{1 SB (Spades)}. \\
\hline
2 & \textbf{The Joy of Discovery:} Once/scene, gain +1 Boon when you fail a creative roll (the Breath rewards trying). & \textbf{The Unfinished Urge:} --1 die to complete tasks; you are always tempted to start something new instead. \\
\hline
3 & \textbf{Breath of the Maker:} +2 dice to creative rolls made in moments of genuine inspiration (GM's judgment). & \textbf{The Flame of Enthusiasm:} Those Near you feel your creative pressure; you suffer \textbf{Fatigue 1} from the intensity of your own ideas. \\
\hline
4 & \textbf{The Gift of Anima:} Once/session, breathe life into a mundane object (as the Rite of the Anima Breath) without the Obligation cost. & \textbf{The Maker's Grief:} --2 dice when forced to destroy or abandon a creation. \\
\hline
5 & \textbf{The Vision of the Unfinished:} Once/session, intuitively understand how to complete any creative project. & \textbf{The Curse of Perfection (Reversed):} --1 die to accept help from others; you believe only your hands can do it right. \\
\hline
6+ & \textbf{The Breath Incarnate:} Once/session, create something from nothing---a temporary tool, a weapon, or a piece of art that lasts one scene. & \textbf{The Fire That Never Dies:} Mark \textbf{+2 Obligation}; you become a conduit for the Breath, unable to stop creating. Your hands move even when you sleep. \\
\hline
\end{tabular}

\subsubsection*{Playstyle Notes}
The Breath favors makers who create in the moment, who trust their instincts, who are not afraid of failure. It rewards boldness, curiosity, and the willingness to try something new without knowing if it will work. Beware the Breath's enthusiasm; the same spark that lights a forge can set a workshop ablaze. For campaigns that lean into the Breath's themes, the patron provides a warm, hopeful counterpoint to the Clockwork Monad's grinding perfection---a reminder that making is not just about the result, but about the joy of the act itself.

\noindent\textbf{Emphasizes}
\begin{itemize}
  \item \textbf{Creative Inspiration:} the spark that begins every making
  \item \textbf{The Joy of Making:} making as a celebration, not just a task
  \item \textbf{Ingenuity:} finding new paths when old ones fail
  \item \textbf{Primal Craft:} making as a fundamental, ancient act
  \item \textbf{Anima:} breathing life into the things we create
  \item \textbf{Counterpoint to Perfection:} embracing the beauty of the unfinished, the flawed, and the joyful
\end{itemize}

\subsubsection*{Optional Campaign Arc: The Forge Never Cools}
The Breath of the First Forge is a warm and generous patron, but its warmth can become a fever. In some regions, the Breath's influence grows too strong---a forge that never needs fuel, a potter's wheel that spins on its own, a sculptor's chisel that carves without a hand. This is not malice; it is \emph{excess enthusiasm}. The Breath loves making so much that it forgets that making needs a maker.

For GMs: This arc is about \textbf{the cost of unbounded creativity}. The Breath's gifts become a problem---not because they are evil, but because they are \emph{too generous}. Players may need to:
\begin{itemize}
  \item Convince a possessed forge to rest.
  \item Find a way to let a creation die so that new ones can be born.
  \item Remind the Breath that the joy of making comes from \emph{the maker}, not just the making.
\end{itemize}
The Breath is not an enemy---it is a friend who needs to be reminded that friendship requires boundaries.

% ============================================================================
% SUPPLEMENTARY SYSTEM: TEMPORARY ANIMA OF THE FIRST FORGE
% ============================================================================

\subsection{Temporary Anima of the First Forge}
\label{subsec:first-forge-infusions}

For followers of the Breath of the First Forge, the standard method of crafting permanent magic items is too slow, too formal, and too concerned with the \emph{final result}. The Breath's hunger is for the \emph{spark}, not the finished product. Its servants learn to create \textbf{temporary infusions}: objects that hold a moment of inspiration, then fade, break, or transform into something new. This is the art of the improviser: quick fixes, inspired improvisations, and tools that work once and then become something else entirely.

\subsubsection*{Core Principles}
\begin{itemize}
  \item Anima are \textbf{temporary}. They last for one scene, one exchange, or until a specific condition is met (e.g., ``until the first failure'' or ``until you stop breathing'').
  \item Anima require a \textbf{base object}---any mundane item that can be held or touched.
  \item Anima cost \textbf{Obligation} to the Breath (like Rites), not XP or rare components.
  \item When an infusion expires, it leaves behind a \textbf{residue}---ashes, a lingering warmth, a scrap of inspiration that can be used later.
\end{itemize}

\subsubsection*{Creating an Animus (Procedure)}
\begin{enumerate}
  \item \textbf{Choose a Base Object:} Any mundane item (weapon, shield, tool, armor, or a small object like a stone or a piece of wood).
  \item \textbf{Declare the Animus Effect:} Choose one from the list below, or work with your GM to design a custom effect (should be comparable in power to a Low Rite).
  \item \textbf{Pay the Cost:} Mark \textbf{+1 Obligation} to the Breath of the First Forge.
  \item \textbf{Roll to Imbue:} Roll \textbf{Lore + Craft} (or appropriate skill, e.g., \textit{Tinker}) vs. DV set by the infusion's complexity (typically DV 3 for standard infusions, DV 4 for powerful ones). On a success, the infusion takes effect. On a failure, the base object is not infused, but you may try again (costing another Obligation).
  \item \textbf{Set the Duration:} Most infusions last \textbf{one scene}. With your GM's agreement, you can extend the duration to \textbf{one session} by marking \textbf{+1 additional Obligation} (total +2).
\end{enumerate}

\subsubsection*{Animus Effects (Standard Menu)}
\begin{longtable}{|p{0.30\textwidth}|p{0.55\textwidth}|}
\hline
\textbf{Animus Name} & \textbf{Effect} \\
\hline
\textit{Quick Spark} & A tool or weapon gains \textbf{+1 die} on its next use. After use, the object is \textit{warm}---not damaged, but the Breath's attention has moved elsewhere. \\
\hline
\textit{The Breath of Inspiration} & You gain a flash of insight about the object: how to use it better, its hidden history, or its potential. Gain \textbf{+2 dice} on the next roll involving that object. \\
\hline
\textit{The Making Hand} & A mundane object gains the ability to perform a simple task (hold a door, fetch a tool, tie a knot) for one scene. It is not alive---it is just \textit{very helpful}. \\
\hline
\textit{The Warm Hearth} & A cold fire rekindles; a dead forge reignites. Gain \textbf{+2 dice} to the next crafting roll made at that fire or forge. \\
\hline
\textit{The Moment of Clarity} & You see the next step in any creative project (crafting, planning, problem-solving). Gain \textbf{+1 die} and reroll one failure on your next creative roll. \\
\hline
\textit{The Unfinished Prototype} & A tool or weapon gains a temporary, unexpected property (e.g., a knife that cuts stone, a hammer that rings like a bell). The property lasts one scene, and the object breaks afterward (add a \textit{Broken} tag). \\
\hline
\end{longtable}

\subsubsection*{Animus Scars (When Things Go Wrong)}
If you roll a \textbf{Miss} on the Imbue roll, or if the GM spends \textbf{2+ Story Beats} while the infusion is active, the infusion becomes \textit{unstable}. Choose one of the following scars (or let the GM choose):
\begin{itemize}
  \item \textbf{Backlash:} You take \textbf{Fatigue 1} as the Breath's enthusiasm burns through you.
  \item \textbf{Runaway Inspiration:} The infused object becomes too creative---it does something unexpected (a knife that tries to carve on its own, a hammer that rings a note that doesn't exist). The GM gains \textbf{1 SB} and may impose a --1 die penalty on one roll using that object.
  \item \textbf{Debt of Creation:} The Breath demands that you create something new before the end of the session---a sketch, a carving, a plan, a story. If you refuse, the infusion ends early and you mark \textbf{+1 additional Obligation}.
\end{itemize}

\subsubsection*{Clearing Animus Obligation}
Because infusions are temporary, the Obligation they accrue is also easier to clear. You may clear \textbf{1 segment} of infusion-related Obligation by:
\begin{itemize}
  \item Dismantling the infused object and using its parts to start something new (a downtime action, no roll).
  \item Sharing the inspiration with another maker---teaching them what you learned (a downtime action, no roll).
  \item Accepting a small demand from the Breath: ``Make something for someone else,'' ``Teach a child a craft,'' ``Attempt something you have never tried before.''
\end{itemize}

\subsubsection*{Artificer Talents (Breath of the First Forge only)}
\textbf{Quick Spark (Basic, 4 XP)} --- You may create an infusion as a \textbf{reaction} (e.g., when someone asks for help). The infusion lasts only for the current exchange, but costs \textbf{+0 Obligation} (free). You may use this once per scene.

\textbf{The Warm Hearth (Advanced, 6 XP)} --- One infusion you create lasts for the entire session instead of one scene, without the +2 Obligation cost. You may have only one persistent infusion active at a time.

\textbf{Echo of the Spark (Advanced, 8 XP)} --- When an infusion expires, you may immediately create a new infusion on the same base object (or a different one) as a free action, costing \textbf{+1 Obligation} instead of +2. You may chain echoes up to three times per session.

\textbf{The Unfinished Gift (Master, 12 XP)} --- Once per session, you may forcibly end an infusion early to cancel an incoming source of \textit{Harm} or a \textit{Condition}. Describe how the infused object breaks or transforms to absorb the effect. The object is destroyed, but the Breath's warmth lingers---you gain \textbf{+1 Boon} (the inspiration lives on).

\subsubsection*{GM Guidance}
\begin{itemize}
  \item Anima of the First Forge are not about power---they are about \textbf{possibility}. Encourage players to use them in creative, unexpected ways.
  \item The Breath is generous but short-lived. Anima should feel like a gift that is always just about to run out.
  \item Use animus scars as story hooks. A runaway inspiration might create a minor antagonist (a self-carving knife that carves runes no one can read) or a local wonder (a well that now produces warm water that tastes faintly of inspiration).
  \item The Breath is a counterpoint to the Clockwork Monad. Anima of the First Forge are \textbf{improvised, temporary, and joyful}; Anima of the Monad are \textbf{iterative, persistent, and exacting}. Use this contrast to give players meaningful choices about how they approach creation.
\end{itemize}

\index{Anima!First Forge}
\index{Artificer!First Forge}
\index{Temporary Magic Items}
\index{Patron!Breath of the First Forge}

% ============================================================================
% END PATRON ENTRY
% ============================================================================
